\section{Introduction and Goal}
\label{sec:intro}

Cryptographic algorithms such as AES-128~\cite{nist2001aes} are mathematically
secure under standard threat models; no practical algebraic attacks are known
against the full cipher.
However, physical implementations of these algorithms routinely leak sensitive
information through observable side channels---most notably through their power
consumption.
Kocher et al.\ demonstrated in 1999~\cite{kocher1999dpa} that an adversary with
physical access to a device can statistically extract a secret key by observing
the power trace of many encryptions; this technique, Differential Power
Analysis (DPA), exploits the fact that CMOS logic gates consume power
proportional to the data they process.

Correlation Power Analysis (CPA), introduced by Brier et al.\ in
2004~\cite{brier2004cpa}, formalises this observation using a Pearson
correlation distinguisher and an explicit leakage model (typically Hamming
Weight of an intermediate value), achieving significantly better efficiency
than DPA in terms of trace complexity.

\subsection*{Project Goal}

The goal of this project is to perform a \emph{quantitative comparative
analysis} of DPA and CPA on a real embedded target---the Raspberry Pi Pico
RP2040 microcontroller~\cite{rp2040datasheet}---running unprotected AES-128.
Specifically, the project targets the following measurements:

\begin{itemize}
  \item \textbf{Key recovery success}: Does each attack recover the full
        128-bit key from the reference trace set?
  \item \textbf{Rank of the correct hypothesis}: What is the rank of the
        correct key byte across all 256 candidates as a function of the
        number of traces used?
  \item \textbf{Leakage model sensitivity}: Does the HW model suffice, or
        do certain bytes require a Hamming-Distance (HD) variant?
  \item \textbf{Trace complexity} ($N_\text{traces}$): How many traces are
        needed before each attack converges (Guessing Entropy reaches 1)?
\end{itemize}

\subsection*{Report Structure}

Section~\ref{sec:threat} defines the threat model and system application.
Section~\ref{sec:vuln} analyses the vulnerability and countermeasure context.
Section~\ref{sec:method} describes the implementation---firmware, acquisition
pipeline, and attack algorithms.
Section~\ref{sec:results} presents experimental results with diagnostic plots.
Section~\ref{sec:discussion} compares DPA and CPA and discusses limitations.
Section~\ref{sec:conclusion} concludes and outlines future work.
Hardware setup details and video evidence are in Appendix~\ref{app:hw};
the software repository structure is in Appendix~\ref{app:repo}.
